\chapter{Normed Spaces and Banach Spaces}

\section{Vector spaces, Hamel bases, and norms}

\begin{notetheorem}[Existence and cardinality of Hamel bases]
Let $X\ne\{0\}$ be a vector space over a field $\mathbb K$.
\begin{enumerate}
  \item $X$ has a Hamel basis.
  \item Any two Hamel bases of $X$ have the same cardinality.
\end{enumerate}
\end{notetheorem}

\begin{proof}
For existence, order the linearly independent subsets of $X$ by inclusion.
The union of a chain is linearly independent, so Zorn's lemma gives a maximal
linearly independent set $E$.  Maximality implies $\Span E=X$.

For uniqueness of cardinality, let $E$ and $F$ be bases.  Every element of
$E$ is a finite linear combination of elements of $F$.  The infinite
Steinitz exchange theorem, equivalently the standard basis-cardinality
argument, gives an injection $E\hookrightarrow F$.  Interchanging $E$ and
$F$ gives an injection $F\hookrightarrow E$.  The Cantor--Bernstein theorem
then yields $|E|=|F|$.
\end{proof}

\begin{noteexercise}[Additive functions on $\mathbb R$]
Let $\mathcal A$ be the set of functions $f\colon\mathbb R\to\mathbb R$
satisfying Cauchy's equation
\[
  f(x+y)=f(x)+f(y).
\]
Describe $\mathcal A$ and determine its cardinality.
\end{noteexercise}

\begin{proof}
Such a function is exactly a $\mathbb Q$-linear map on the
$\mathbb Q$-vector space $\mathbb R$.  Fix a Hamel basis $H$ of
$\mathbb R$ over $\mathbb Q$.  An arbitrary function $a\colon H\to\mathbb R$
extends uniquely to an additive function by
\[
  f\!\left(\sum_{j=1}^n q_jh_j\right)=\sum_{j=1}^nq_ja(h_j).
\]
Conversely every additive function is obtained this way.  Since
$|H|=|\mathbb R|=\mathfrak c$, one has
\[
  |\mathcal A|=|\mathbb R|^{\mathfrak c}=2^{\mathfrak c}.
\]
Only the choices $a(h)=ch$ give continuous additive functions.
\end{proof}

\begin{notelemma}[Extension of a linearly independent family]
Every linearly independent family in a vector space is contained in a Hamel
basis.
\end{notelemma}

\begin{proof}
Apply Zorn's lemma to the linearly independent sets containing the given
family.  A maximal member spans the whole space.
\end{proof}

\begin{notedefinition}[Equivalent norms]
Two norms $\|\cdot\|_1$ and $\|\cdot\|_2$ on a vector space $X$ are
\emph{equivalent} if there are constants $c,C>0$ such that
\[
  c\|x\|_1\le\|x\|_2\le C\|x\|_1
  \qquad(x\in X).
\]
\end{notedefinition}

\begin{notetheorem}[Equivalent norms and topology]
Two norms on $X$ are equivalent if and only if they induce the same topology.
\end{notetheorem}

\begin{proof}
If the topologies agree, the $\|\cdot\|_2$-unit ball is a neighbourhood of
$0$ for $\|\cdot\|_1$, so some $r>0$ satisfies
$B_1(0,r)\subseteq B_2(0,1)$.  Applying this inclusion to
$rx/(2\|x\|_1)$ gives
$\|x\|_2\le 2r^{-1}\|x\|_1$ for $x\ne0$.  Reversing the roles of the
norms gives the other inequality.

Conversely, the displayed two-sided estimate implies
\[
  B_1\!\left(x,\frac rC\right)\subseteq B_2(x,r),
  \qquad
  B_2(x,cr)\subseteq B_1(x,r).
\]
Thus the open sets for either norm are open for the other.
\end{proof}

\begin{notetheorem}[Separable spaces are increasing unions]
Let $X$ be a separable infinite-dimensional normed space.  Then there are
finite-dimensional subspaces
\[
  X_1\subsetneq X_2\subsetneq\cdots,
  \qquad \dim X_n=n,
\]
such that $\bigcup_nX_n$ is dense in $X$.
\end{notetheorem}

\begin{proof}
Let $(d_m)$ be dense.  Scan this sequence in order and retain $d_m$ exactly
when it is not in the span of the previously retained vectors.  Denote the
resulting linearly independent sequence by $(x_n)$ and put
$X_n=\Span\{x_1,\ldots,x_n\}$.  Every discarded $d_m$ already belongs to
one of these spans, so $\bigcup_nX_n$ contains a dense subset.  Since $X$ is
infinite-dimensional, the retained sequence is infinite, and
$\dim X_n=n$.
\end{proof}

\begin{notetheorem}[Every vector space can be normed]
Every vector space admits a norm.
\end{notetheorem}

\begin{proof}
Choose a Hamel basis $(e_i)_{i\in I}$.  Every $x\in X$ has a unique finite
expansion $x=\sum_{i\in F}a_ie_i$.  Define
\[
  \|x\|=\sum_{i\in F}|a_i|.
\]
Uniqueness of the expansion makes this well defined, and the norm axioms are
immediate.
\end{proof}

\section{Distance to compact sets and spaces of functions}

\begin{noteproposition}[Distance to a compact set]
Let $K$ be a nonempty compact subset of a normed space $X$ and define
\[
  d_K(x)=\inf_{z\in K}\|x-z\|.
\]
Then the infimum is attained for every $x$, and
\[
  |d_K(x)-d_K(y)|\le\|x-y\|.
\]
In particular, $d_K$ is continuous.
\end{noteproposition}

\begin{proof}
Choose $z_n\in K$ with $\|x-z_n\|\to d_K(x)$.  Compactness gives a
convergent subsequence $z_{n_k}\to z\in K$, and continuity of the norm gives
$\|x-z\|=d_K(x)$.  For any $z\in K$,
\[
  d_K(x)\le\|x-z\|\le\|x-y\|+\|y-z\|.
\]
Taking the infimum over $z$ gives $d_K(x)-d_K(y)\le\|x-y\|$; interchange
$x$ and $y$.
\end{proof}

\begin{noteremark}[Nearest-point selections]
A nearest point need not be unique, and an arbitrary selection
$P(x)\in K$ with $\|x-P(x)\|=d_K(x)$ need not be continuous.  For example,
for $K=\{-1,1\}\subset\mathbb R$, any nearest-point selection jumps at
$0$.  What is always continuous is the distance function $d_K$.
\end{noteremark}

\begin{notetheorem}[Supremum norm on continuous maps]
Let $K$ be compact and $Y$ a normed space.  The formula
\[
  \|f\|_\infty=\sup_{x\in K}\|f(x)\|
\]
defines a norm on $C(K;Y)$.
\end{notetheorem}

\begin{notetheorem}[Bounded maps]
For arbitrary $X$ and a normed space $Y$, let $B(X;Y)$ be the vector space of
bounded maps $f\colon X\to Y$.  Then $\|f\|_\infty$ is a norm on
$B(X;Y)$.
\end{notetheorem}

\begin{notetheorem}[Uniform limits preserve continuity]
If $f_n\colon X\to Y$ are continuous and $f_n\to f$ uniformly, then $f$ is
continuous.
\end{notetheorem}

\begin{proof}
Fix $x\in X$ and $\varepsilon>0$.  Choose $n$ with
$\|f-f_n\|_\infty<\varepsilon/3$, then choose a neighbourhood $U$ of $x$
so that $\|f_n(y)-f_n(x)\|<\varepsilon/3$ for $y\in U$.  The triangle
inequality gives $\|f(y)-f(x)\|<\varepsilon$.
\end{proof}

\begin{notecorollary}
The subspace $B(X;Y)\cap C(X;Y)$ is closed in $B(X;Y)$ for the supremum
norm.
\end{notecorollary}

\section{Banach spaces and completeness}

\begin{notedefinition}[Banach space]
A normed space is a \emph{Banach space} if it is complete for the metric
$d(x,y)=\|x-y\|$.
\end{notedefinition}

\begin{notedefinition}[Schauder basis]
A sequence $(e_n)$ in a normed space $X$ is a \emph{Schauder basis} if every
$x\in X$ has a unique expansion
\[
  x=\sum_{n=1}^{\infty}a_ne_n
\]
converging in norm.  Unlike a Hamel basis, a Schauder basis uses infinite
series and is therefore a topological notion.
\end{notedefinition}

\begin{noteremark}
A normed space with a Schauder basis is separable, but not every separable
Banach space has a Schauder basis.
\end{noteremark}

\begin{notelemma}[Translation and scalar invariance]
The metric induced by a norm satisfies
\[
  d(x+a,y+a)=d(x,y),
  \qquad
  d(\alpha x,\alpha y)=|\alpha|d(x,y).
\]
\end{notelemma}

\begin{proof}
Both identities follow immediately from
$\|(x+a)-(y+a)\|=\|x-y\|$ and
$\|\alpha(x-y)\|=|\alpha|\|x-y\|$.
\end{proof}

\begin{notetheorem}[Closed subspaces of Banach spaces]
Let $X$ be a Banach space and $Y\subseteq X$ a linear subspace.  Then $Y$ is
Banach with the inherited norm if and only if $Y$ is closed in $X$.
\end{notetheorem}

\begin{proof}
If $Y$ is closed, a Cauchy sequence in $Y$ converges in $X$ and its limit lies
in $Y$.  Conversely, if $Y$ is complete and $y_n\in Y$ converges in $X$ to
$x$, then $(y_n)$ is Cauchy in $Y$ and converges there to some $y\in Y$.
Uniqueness of limits gives $x=y$.
\end{proof}

\begin{notedefinition}[Series in normed spaces]
The series $\sum_{n=1}^{\infty}x_n$ converges if its partial sums converge.
It is \emph{absolutely convergent} if $\sum_n\|x_n\|<\infty$.
\end{notedefinition}

\begin{notetheorem}[Isomorphic images of Banach spaces]
Let $X$ be Banach, let $Y$ be normed, and let
$A\colon X\to Y$ be a bounded bijection with bounded inverse.  Then $Y$ is
Banach.
\end{notetheorem}

\begin{proof}
If $(y_n)$ is Cauchy in $Y$, then
$\|A^{-1}y_n-A^{-1}y_m\|\le\|A^{-1}\|\|y_n-y_m\|$, so
$A^{-1}y_n\to x\in X$.  Hence $y_n\to Ax$.
\end{proof}

\begin{notetheorem}[Closed range from a lower bound]
Let $X$ be Banach, let $Y$ be normed, and let $A\in\Bop{X}{Y}$.  If
\[
  \|Ax\|\ge c\|x\|
  \qquad(x\in X)
\]
for some $c>0$, then $A$ is injective and $\Ran(A)$ is closed in $Y$.
\end{notetheorem}

\begin{proof}
Injectivity is immediate.  If $Ax_n\to y$, then
$\|x_n-x_m\|\le c^{-1}\|Ax_n-Ax_m\|$, so $(x_n)$ is Cauchy and converges
to some $x\in X$.  Continuity gives $Ax=y$.
\end{proof}

\begin{notetheorem}[Neumann series]
Let $X$ be Banach and $A\in\Bop{X}{X}$ with $\|A\|<1$.  Then
\[
  (I-A)^{-1}=\sum_{n=0}^{\infty}A^n
\]
in operator norm, and
\[
  \|(I-A)^{-1}\|\le\frac1{1-\|A\|}.
\]
\end{notetheorem}

\begin{proof}
The series converges absolutely in the Banach algebra $\Bop{X}{X}$.  Its
partial sums $S_N$ satisfy
$(I-A)S_N=S_N(I-A)=I-A^{N+1}$.  Letting $N\to\infty$ gives the inverse.
\end{proof}

\begin{notecorollary}[The invertible group is open]
If $A\in\Bop{X}{X}$ is invertible and
$\|B-A\|<\|A^{-1}\|^{-1}$, then $B$ is invertible and
\[
  \|B^{-1}-A^{-1}\|
  \le
  \frac{\|A^{-1}\|^2\|B-A\|}
       {1-\|A^{-1}\|\|B-A\|}.
\]
Thus inversion is continuous on the group of invertible operators.
\end{notecorollary}

\begin{notetheorem}[A Schauder basis implies separability]
If a normed space has a Schauder basis, then it is separable.
\end{notetheorem}

\begin{proof}
Let $(e_n)$ be a Schauder basis.  The set of finite rational (or Gaussian
rational) linear combinations of the $e_n$ is countable.  Approximation of
the coefficients in a sufficiently long partial sum shows this set is dense.
\end{proof}

\begin{noteexercise}[A power-small operator]
Let $X$ be Banach and suppose $A\in\Bop{X}{X}$ satisfies $\|A^p\|<1$ for
some $p\ge1$.  Show that $I-A$ is invertible.
\end{noteexercise}

\begin{proof}
Since
\[
  I-A^p=(I-A)(I+A+\cdots+A^{p-1}),
\]
the Neumann theorem makes $I-A^p$ invertible.  The factors commute, so
\[
  (I-A)^{-1}=(I+A+\cdots+A^{p-1})(I-A^p)^{-1}.
\]
\end{proof}

\begin{notetheorem}[Absolute convergence characterizes Banach spaces]
A normed space $X$ is Banach if and only if every series satisfying
$\sum_n\|x_n\|<\infty$ converges in $X$.
\end{notetheorem}

\begin{proof}
In a Banach space the partial sums are Cauchy because tails have norm bounded
by the corresponding scalar tail.

Conversely, let $(x_n)$ be Cauchy.  Choose a subsequence $(x_{n_k})$ with
\[
  \|x_{n_{k+1}}-x_{n_k}\|<2^{-k}.
\]
The series
$x_{n_1}+\sum_k(x_{n_{k+1}}-x_{n_k})$ converges by hypothesis, so the
subsequence converges.  A Cauchy sequence with a convergent subsequence
converges to the same limit.
\end{proof}

\begin{notetheorem}[Completion of a normed space]
Every normed space $X$ has a Banach completion $\widehat X$ in which $X$ is
isometrically embedded as a dense subspace.  It is unique up to an isometric
isomorphism fixing $X$.
\end{notetheorem}

\begin{notedefinition}[Seminorm]
A map $p\colon X\to[0,\infty)$ is a seminorm if
\[
  p(\alpha x)=|\alpha|p(x),
  \qquad
  p(x+y)\le p(x)+p(y).
\]
Then $p(0)=0$ and $|p(x)-p(y)|\le p(x-y)$.
\end{notedefinition}

\begin{noteproposition}[Quotient by the null space]
Let $N=\{x:p(x)=0\}$.  Then $N$ is a subspace and
\[
  \|x+N\|_{X/N}=p(x)
\]
is a well-defined norm on $X/N$.
\end{noteproposition}

\begin{proof}
If $x-y\in N$, then
$|p(x)-p(y)|\le p(x-y)=0$, so the formula is independent of the chosen
representative.  Positive definiteness on the quotient and the remaining norm
axioms follow directly.
\end{proof}

\section{Finite-dimensional normed spaces}

\begin{notelemma}[Linear combinations]
If $x_1,\ldots,x_n$ are linearly independent in a normed space, then there is
$c>0$ such that
\[
  \left\|\sum_{j=1}^n a_jx_j\right\|
  \ge c\sum_{j=1}^n|a_j|
  \qquad(a_1,\ldots,a_n\in\mathbb K).
\]
\end{notelemma}

\begin{proof}
On the compact set
$S=\{a\in\mathbb K^n:\sum_j|a_j|=1\}$, the continuous function
$a\mapsto\|\sum_ja_jx_j\|$ has a minimum.  The minimum is positive, because
zero would give a nontrivial linear dependence.  Homogeneity yields the
estimate.
\end{proof}

\begin{notetheorem}[Completeness]
Every finite-dimensional normed space is Banach.  More generally, every
finite-dimensional subspace of a normed space is complete.
\end{notetheorem}

\begin{proof}
Write $y_m=\sum_{j=1}^na_j^{(m)}e_j$ in a basis.  The preceding lemma shows
that a Cauchy sequence $(y_m)$ has each scalar coordinate
$(a_j^{(m)})$ Cauchy.  The coordinates converge in $\mathbb K$, and the
corresponding vector is the limit of $(y_m)$.
\end{proof}

\begin{notecorollary}
Every finite-dimensional subspace of a normed space is closed.
\end{notecorollary}

\begin{notetheorem}[Equivalence of finite-dimensional norms]
Any two norms on a finite-dimensional vector space are equivalent.
\end{notetheorem}

\begin{proof}
Choose a basis and compare each norm with the coordinate $\ell^1$ norm.  The
upper bounds follow from the triangle inequality, and the lower bounds follow
from the linear-combination lemma.
\end{proof}

\begin{notetheorem}[Inequivalent norms in infinite dimension]
Every infinite-dimensional vector space admits two inequivalent norms.
\end{notetheorem}

\begin{proof}
Choose a Hamel basis containing a countable subfamily $(e_n)$.  For a finite
expansion $x=\sum a_ie_i$, set
\[
  \|x\|_1=\sum_i|a_i|,
  \qquad
  \|x\|_2=\sum_iw_i|a_i|,
\]
where $w_n=n$ on the chosen subfamily and $w_i=1$ elsewhere.  Then
$\|e_n\|_1=1$ but $\|e_n\|_2=n$, so no uniform comparison
$\|x\|_2\le C\|x\|_1$ is possible.
\end{proof}

\section{Compactness}

\begin{notedefinition}[Sequential compactness]
A subset $M$ of a metric space is compact if every sequence in $M$ has a
subsequence converging to an element of $M$.  In metric spaces this is
equivalent to the open-cover definition.
\end{notedefinition}

\begin{notelemma}
Every compact subset of a metric space is closed and bounded.
\end{notelemma}

\begin{proof}
If $M$ were unbounded, one could choose $x_n\in M$ with
$d(x_n,x_0)>n$, and no subsequence could converge.  If a sequence in $M$
converges in the ambient space, compactness gives a subsequence converging to
a point of $M$; uniqueness of the limit puts the original limit in $M$.
\end{proof}

\begin{noteexample}[The converse fails]
In $\ell^2$, the set $\{e_n:n\ge1\}$ of standard unit vectors is closed and
bounded, but not compact, because
\[
  \|e_n-e_m\|_2=\sqrt2\qquad(n\ne m).
\]
Thus it has no convergent subsequence.
\end{noteexample}

\begin{notetheorem}[Heine--Borel in finite dimension]
A subset of a finite-dimensional normed space is compact if and only if it is
closed and bounded.
\end{notetheorem}

\begin{proof}
Use a basis to identify the space linearly and homeomorphically with
$\mathbb K^n$.  Boundedness of the vectors implies boundedness of their
coordinate tuples; Bolzano--Weierstrass gives a convergent coordinate
subsequence, and closedness keeps the limit in the set.
\end{proof}

\begin{notelemma}[Riesz's lemma]
Let $Y$ be a proper closed subspace of a normed space $X$.  For every
$0<\theta<1$ there is $z\in X$ such that
\[
  \|z\|=1,
  \qquad
  \|z-y\|>\theta\quad(y\in Y).
\]
\end{notelemma}

\begin{proof}
Choose $x\notin Y$ and put $d=\dist(x,Y)>0$.  Select $y_0\in Y$ with
$\|x-y_0\|<d/\theta$ and set
$z=(x-y_0)/\|x-y_0\|$.  For $y\in Y$,
\[
  \|z-y\|
  =\frac{\|x-(y_0+\|x-y_0\|y)\|}{\|x-y_0\|}
  \ge\frac d{\|x-y_0\|}>\theta.
\]
\end{proof}

\begin{noteremark}
The strict endpoint $\theta=1$ is impossible in the displayed form, because
$y=0$ gives $\|z-y\|=1$.  If a nearest point to $x$ in $Y$ exists, one may
obtain the sharp non-strict statement
$\dist(z,Y)=1$, equivalently $\|z-y\|\ge1$ for every $y\in Y$.
\end{noteremark}

\begin{notetheorem}[Compact unit ball forces finite dimension]
If the closed unit ball of a normed space $X$ is compact, then $X$ is
finite-dimensional.
\end{notetheorem}

\begin{proof}
If $X$ were infinite-dimensional, repeatedly apply Riesz's lemma to
$Y_n=\Span\{x_1,\ldots,x_n\}$ to obtain unit vectors with
$\|x_n-x_m\|>1/2$ for $n\ne m$.  This bounded sequence has no convergent
subsequence, contradicting compactness.
\end{proof}

\begin{noteremark}
A metric space is compact if and only if it is complete and totally bounded.
\end{noteremark}

\begin{notetheorem}[Banach space of continuous functions]
Let $K$ be compact and $Y$ Banach.  Then $C(K;Y)$ is Banach under the
supremum norm.
\end{notetheorem}

\begin{proof}
If $(f_n)$ is Cauchy in $\|\cdot\|_\infty$, then for each $x$ the sequence
$(f_n(x))$ converges in $Y$.  Let $f(x)$ be the pointwise limit.  Passing to
the limit in the uniform Cauchy estimate shows $f_n\to f$ uniformly, and the
uniform-limit theorem makes $f$ continuous.
\end{proof}

\begin{notecorollary}
If $Y$ is Banach, then $B(X;Y)$ is Banach under the supremum norm.
\end{notecorollary}

\begin{notetheorem}[Best approximation in finite-dimensional subspaces]
Let $Y$ be a finite-dimensional subspace of a normed space $X$.  For every
$x\in X$ there is $y_0\in Y$ such that
\[
  \|x-y_0\|=\dist(x,Y).
\]
If, in addition, $X$ is strictly convex, then $y_0$ is unique.
\end{notetheorem}

\begin{proof}
Choose a minimizing sequence in $Y$.  It is bounded, so finite-dimensional
compactness gives a convergent subsequence; continuity of the norm gives a
minimizer.  If $y_1\ne y_2$ are minimizers and $X$ is strictly convex, then
\[
 \left\|x-\frac{y_1+y_2}{2}\right\|
 =\left\|\frac{(x-y_1)+(x-y_2)}2\right\|
 <\dist(x,Y),
\]
a contradiction.
\end{proof}

\begin{noteexample}
The spaces $\ell^p$ are strictly convex for $1<p<\infty$, but not for
$p=1$ or $p=\infty$.  For example, in $\ell^1_2$, let
$Y=\Span\{(1,1)\}$ and $x=(1,0)$.  Then
\[
  \|x-t(1,1)\|_1=|1-t|+|t|=1
  \qquad(0\le t\le1),
\]
so the best approximant is not unique.
\end{noteexample}

\begin{notetheorem}
Every compact subset of an infinite-dimensional normed space has empty
interior.
\end{notetheorem}

\begin{proof}
If a compact set contained a ball $B(x,r)$, then its closed subset
$\overline B(x,r/2)$ would be compact.  Translation and scaling would make
the closed unit ball compact, contradicting the preceding theorem.
\end{proof}

\section{Linear and bounded operators}

\begin{notedefinition}[Inverse operator]
A linear map $T\colon X\to Y$ is invertible if it is bijective.  Its inverse
$T^{-1}\colon Y\to X$ is then linear.  If $S\colon X\to Y$ and
$T\colon Y\to Z$ are bijective, then
\[
  (TS)^{-1}=S^{-1}T^{-1}.
\]
\end{notedefinition}

\begin{notedefinition}[Bounded linear map]
For normed spaces $V,W$ and a linear map $T\colon V\to W$, define
\[
  \|T\|=\sup_{\|v\|\le1}\|Tv\|.
\]
The map is \emph{bounded} if $\|T\|<\infty$.  The vector space of bounded
linear maps is denoted $\Bop{V}{W}$.
\end{notedefinition}

\begin{notetheorem}[Bounded inverse and a lower bound]
Let $A\colon X\to Y$ be a linear map between normed spaces.  The following
are equivalent.
\begin{enumerate}
  \item $A$ is injective and the inverse
        $A^{-1}\colon\Ran(A)\to X$ is bounded.
  \item There is $c>0$ such that
        $\|Ax\|\ge c\|x\|$ for every $x\in X$.
\end{enumerate}
In that case the best lower-bound constant is $c=\|A^{-1}\|^{-1}$.
\end{notetheorem}

\begin{proof}
If $A^{-1}$ is bounded, then
$\|x\|\le\|A^{-1}\|\|Ax\|$.  Conversely, the lower bound gives
injectivity and, for $y=Ax$,
$\|A^{-1}y\|=\|x\|\le c^{-1}\|y\|$.
\end{proof}

\begin{noteexercise}[A functional on $c_0$]
Let $(\alpha_n)\in\ell^1$ and define
\[
  f(x)=\sum_{n=1}^{\infty}\alpha_nx_n,
  \qquad x=(x_n)\in c_0.
\]
Then $f\in c_0^*$ and
\[
  \|f\|=\sum_{n=1}^{\infty}|\alpha_n|.
\]
If infinitely many $\alpha_n$ are nonzero, the norm need not be attained on
the unit sphere of $c_0$.
\end{noteexercise}

\begin{proof}
The estimate $|f(x)|\le\|x\|_\infty\sum|\alpha_n|$ gives the upper bound.
For $N\ge1$, choose $x^{(N)}$ supported on the first $N$ coordinates with
$x_n^{(N)}$ equal to the phase of $\alpha_n$.  Then
$\|x^{(N)}\|_\infty=1$ and
$|f(x^{(N)})|=\sum_{n\le N}|\alpha_n|$, proving equality after
$N\to\infty$.  Equality for one $x\in c_0$ would require
$|x_n|=1$ with the correct phase whenever $\alpha_n\ne0$, which is impossible
when infinitely many coefficients are nonzero.
\end{proof}

\begin{notetheorem}[Operator norm]
The operator norm makes $\Bop{V}{W}$ a normed space.  It satisfies
\[
  \|ST\|\le\|S\|\,\|T\|,
  \qquad
  \|Tv\|\le\|T\|\,\|v\|.
\]
If $W$ is Banach, then $\Bop{V}{W}$ is Banach.
\end{notetheorem}

\begin{proof}
Only completeness needs comment.  If $(T_n)$ is Cauchy in operator norm, then
$(T_nv)$ is Cauchy in $W$ for each $v$.  Define $Tv=\lim_nT_nv$.
Linearity passes to the limit, the uniform operator bound gives boundedness,
and the Cauchy estimate passes to the limit to show $\|T_n-T\|\to0$.
\end{proof}

\begin{notetheorem}[Finite-dimensional domain]
Every linear map from a finite-dimensional normed space into a normed space is
bounded.
\end{notetheorem}

\begin{proof}
For a basis $(e_j)_{j=1}^n$, write $x=\sum a_je_j$.  The
linear-combination lemma bounds $\sum|a_j|$ by a constant multiple of
$\|x\|$, and therefore
\[
  \|Tx\|\le\sum_j|a_j|\|Te_j\|\le C\|x\|.
\]
\end{proof}

\begin{notetheorem}[Characterization by linear functionals]
A normed space is finite-dimensional if and only if every linear functional
on it is continuous.
\end{notetheorem}

\begin{proof}
The forward direction is the preceding theorem.  If $X$ is
infinite-dimensional, choose linearly independent unit vectors $(e_n)$ and
extend them to a Hamel basis.  Define a linear functional by
$\varphi(e_n)=n$ and by zero on the other basis vectors.  Then
$\|e_n\|=1$ but $|\varphi(e_n)|=n$, so $\varphi$ is unbounded.
\end{proof}

\begin{notetheorem}[Continuity is equivalent to boundedness]
A linear map between normed spaces is continuous if and only if it is bounded.
It is enough that it be continuous at one point.
\end{notetheorem}

\begin{proof}
A bounded map is Lipschitz.  Conversely, continuity at $0$ gives
$\delta>0$ such that $\|x\|<\delta$ implies $\|Tx\|<1$.
For $x\ne0$, apply this to $(\delta/2\|x\|)x$ to obtain
$\|Tx\|\le2\delta^{-1}\|x\|$.  Continuity at an arbitrary point translates
to continuity at $0$ by linearity.
\end{proof}

\begin{notecorollary}
For a bounded linear operator $T$, the null space $\Ker(T)$ is closed.
\end{notecorollary}

\begin{notedefinition}[Compact operator]
A linear map $A\colon X\to Y$ is \emph{compact} if it maps every bounded
sequence in $X$ to a sequence with a convergent subsequence in $Y$.
Equivalently, the image of every bounded set is relatively compact.  A compact
linear map is automatically bounded.
\end{notedefinition}

\begin{notetheorem}[Basic compact-operator criteria]
Let $A\colon X\to Y$ be linear.
\begin{enumerate}
  \item Every compact operator is bounded.
  \item $A$ is compact if and only if every bounded sequence $(x_n)$ has a
        subsequence $(x_{n_k})$ for which $(Ax_{n_k})$ converges.
  \item Every bounded finite-rank operator is compact.
\end{enumerate}
\end{notetheorem}

\begin{proof}
If the image of the unit ball were unbounded, one could choose a sequence whose
image has no convergent subsequence.  The sequential criterion is exactly the
definition of relative compactness in metric spaces.  For finite rank, the
image of a bounded set is bounded in a finite-dimensional space, whose closure
is compact.
\end{proof}

\begin{noteexercise}
The inverse of a bounded bijective operator need not be bounded when the
codomain is incomplete.  Let
\[
  T\colon \ell^1\longrightarrow \Ran(T)\subset\ell^1,
  \qquad
  T(x_n)=\left(\frac{x_n}{n}\right),
\]
where the range carries the inherited $\ell^1$ norm.  Then $T$ is a bounded
bijection onto its range, but
$\|T^{-1}e_n/n\|_1=1$ while $\|e_n/n\|_1=1/n$, so $T^{-1}$ is
unbounded.  Accordingly, $\Ran(T)$ is not complete.
\end{noteexercise}

\begin{notetheorem}[Extension from a dense subspace]
Let $U$ be a dense subspace of a normed space $V$, let $W$ be Banach, and let
$S\in\Bop{U}{W}$.  Then there is a unique
$T\in\Bop{V}{W}$ extending $S$, and
\[
  \|T\|=\|S\|.
\]
\end{notetheorem}

\begin{proof}
For $v\in V$, choose $u_n\in U$ with $u_n\to v$ and define
$Tv=\lim_nSu_n$.  The sequence $(Su_n)$ is Cauchy, and completeness of $W$
provides the limit.  The estimate
$\|Su_n-Su_n'\|\le\|S\|\|u_n-u_n'\|$ shows independence of the
approximating sequence.  Passing to limits gives linearity and
$\|Tv\|\le\|S\|\|v\|$; restriction gives the reverse norm inequality.
Uniqueness follows from density and continuity.
\end{proof}

\begin{noteexample}[Completeness of the codomain is essential]
The identity map from $\mathbb Q$ to $\mathbb Q$, with the usual norm and
viewed on the dense subspace $\mathbb Q\subset\mathbb R$, has no continuous
extension $\mathbb R\to\mathbb Q$.
\end{noteexample}

\begin{notecorollary}[Composition estimate]
If $T\in\Bop{U}{V}$ and $S\in\Bop{V}{W}$, then
\[
  \|S\circ T\|\le\|S\|\,\|T\|.
\]
\end{notecorollary}

\begin{noteexercise}[Subspaces and connectedness]
\begin{enumerate}
  \item A linear subspace containing a nonempty open ball is the whole space.
  \item The only subsets of a normed space that are both open and closed are
        $\varnothing$ and the whole space.
  \item The closure of a linear subspace is again a linear subspace.
\end{enumerate}
\end{noteexercise}

\begin{proof}
For (1), translate and scale the ball to show every vector belongs to the
subspace.  For (2), normed spaces are path connected via line segments.
For (3), take convergent sequences from the subspace and pass addition and
scalar multiplication to the limit.
\end{proof}

\begin{notedefinition}[Quotient norm]
Let $U$ be a subspace of a normed space $V$.  On the algebraic quotient set
$V/U$ define
\[
  \|v+U\|=\inf_{u\in U}\|v+u\|=\dist(v,U).
\]
\end{notedefinition}

\begin{notetheorem}[Quotient spaces]
The quotient formula is a norm if and only if $U$ is closed.  Moreover:
\begin{enumerate}
  \item if $V$ is Banach and $U$ is closed, then $V/U$ is Banach;
  \item if $U$ and $V/U$ are Banach, then $V$ is Banach.
\end{enumerate}
\end{notetheorem}

\begin{proof}
The only nonautomatic norm axiom is positive definiteness:
$\|v+U\|=0$ exactly when $v\in\overline U$.  Thus it holds precisely when
$U$ is closed.

For the first completeness statement, from a summable series in $V/U$ choose
representatives with nearly minimal norm; the resulting corrected series in
$V$ converges and projects to the required quotient sum.  For the converse,
let $\sum v_n$ be absolutely convergent in $V$.  Its quotient series
converges, so after subtracting a suitable $v\in V$ its partial sums approach
$U$.  Choose corrections in $U$ with summable errors; completeness of $U$
then gives convergence in $V$.
\end{proof}

\begin{notetheorem}[Equivalent boundedness criteria]
For a linear map $T\colon V\to W$ between normed spaces, the following are
equivalent:
\begin{enumerate}
  \item $T$ is bounded;
  \item $T$ is continuous at some point;
  \item $T$ is uniformly continuous;
  \item $T^{-1}(B_W(0,1))$ contains a neighbourhood of $0$.
\end{enumerate}
\end{notetheorem}

\begin{proof}
Boundedness gives a global Lipschitz estimate.  Continuity at one point
translates to continuity at zero.  If a ball $B_V(0,\delta)$ is mapped into
$B_W(0,1)$, homogeneity gives $\|Tv\|\le2\delta^{-1}\|v\|$.
\end{proof}

\section{Dual spaces}

\begin{notedefinition}[Algebraic and continuous duals]
The algebraic dual of a vector space $X$ is the vector space of all linear
functionals on $X$.  For a normed space $X$, the \emph{dual space} $X^*$ is
the subspace of bounded linear functionals, with norm
\[
  \|f\|=\sup_{\|x\|\le1}|f(x)|.
\]
\end{notedefinition}

\begin{notetheorem}
The dual $X^*$ of every normed space is Banach.
\end{notetheorem}

\begin{proof}
This is the completeness theorem for $\Bop{X}{\mathbb K}$, since the scalar
field is complete.
\end{proof}

\begin{noteexercise}[Codimension-one kernels]
Let $f\ne0$ be a linear functional on a vector space $X$ and choose
$x_0$ with $f(x_0)\ne0$.
\begin{enumerate}
  \item Every $x\in X$ has a unique representation
        $x=\alpha x_0+y$ with $y\in\Ker f$.
  \item If $f_1,f_2\ne0$ have the same kernel, then $f_1$ is a nonzero scalar
        multiple of $f_2$.
\end{enumerate}
\end{noteexercise}

\begin{proof}
Take $\alpha=f(x)/f(x_0)$ and $y=x-\alpha x_0$.  Uniqueness follows by
applying $f$.  For the second claim, use the same decomposition and compare
the values on a vector outside the common kernel.
\end{proof}
